\unitopen{22}
\begin{candos}
\item napisać krótki, uprzejmy e-mail służbowy z tematem, zwrotem początkowym i zakończeniem;
\item zadzwonić do firmy, przedstawić się i poprosić o połączenie;
\item przełożyć lub potwierdzić termin spotkania;
\item zostawić i przyjąć wiadomość telefoniczną.
\end{candos}

\section{Słownictwo}
\vocabtable{22}

\section{E-mail służbowy}
\begin{center}\small
\begin{tabular}{@{}p{0.22\linewidth}p{0.70\linewidth}@{}}
\toprule
\textsf{\textbf{Element}} & \textsf{\textbf{Wzór}}\\
\midrule
\de{Betreff} & \de{Terminverschiebung: Meeting am Dienstag}\\
Zwrot formalny & \de{Sehr geehrte Frau Keller,} / \de{Sehr geehrter Herr Nowak,}\\
Zwrot półformalny & \de{Guten Tag Frau Keller,} / \de{Liebe Anna,} (po tykaniu)\\
Cel & \de{ich schreibe Ihnen, weil \dots{}} / \de{Leider kann ich am Dienstag nicht teilnehmen.}\\
Prośba & \de{Könnten wir den Termin auf Donnerstag verschieben?}\\
Zakończenie & \de{Vielen Dank für Ihre Rückmeldung.}\\
Pozdrowienie & \de{Mit freundlichen Grüßen} / CH: \de{Freundliche Grüsse}\\
\bottomrule
\end{tabular}
\end{center}
Po zwrocie \emph{Sehr geehrte Frau Keller,} pierwsze słowo pisze się \emph{małą} literą, o ile nie jest rzeczownikiem.
Do kolegów: \de{Hallo Anna,} + \de{Viele Grüße}. Załącznik: \de{Im Anhang finden Sie \dots}

\section{Telefon}
\de{Keller, guten Tag.} (odbierający) \de{Guten Tag, hier spricht Piotr Nowak von der Firma Alpha. Kann ich bitte Frau Keller sprechen?}
\de{Einen Moment, ich verbinde Sie.} \de{Sie ist leider nicht im Haus. Kann ich ihr etwas ausrichten?}
\de{Ich rufe später noch einmal an.} Przeliteruj nazwisko: \de{Das schreibt sich N-O-W-A-K.}

\begin{falle}
Polski e-mail zaczyna się często od „Witam”, co po niemiecku nie ma odpowiednika grzecznościowego: \de{Hallo}
jest dla znajomych, a formalnie piszesz \de{Sehr geehrte \dots}. Drugi błąd: kalka „Z poważaniem” tłumaczona jako
\de{Mit Respekt}; poprawnie \de{Mit freundlichen Grüßen}. Trzeci: nie zaczynaj kolejnego zdania wielką literą po przecinku
po zwrocie początkowym.
\end{falle}

\begin{szwajcaria}
W Szwajcarii pisze się \de{Freundliche Grüsse} (bez \de{ß}), a w firmach także krócej \de{Beste Grüsse} lub \de{Liebe Grüsse} do
znajomych. Formalność jest mniejsza niż w Niemczech: do współpracowników po imieniu, ale e-mail do
klienta czy urzędu zachowaj formalny. \de{Grüezi Frau Keller,} jest także często spotykane na początku maila.
Posiedzenie w CH to \de{die Sitzung}, a ustalanie terminu \de{abmachen}.
\end{szwajcaria}

\begin{praca}
Przykład: \de{Sehr geehrter Herr Nowak, vielen Dank für Ihre Nachricht. Leider kann ich am Dienstag nicht teilnehmen.
Könnten wir den Termin auf Donnerstag, 14 Uhr, verschieben? Im Anhang finden Sie die Unterlagen. Mit freundlichen Grüßen
Anna Keller.} Zapisz go jako własny szablon (3 warianty: formalny, półformalny, zespół).
\end{praca}

\section{Dialogi}
\textbf{Dialog 1 --- telefon do firmy}
\begin{dialog}
\turn{A}{Firma Alpha, Keller, guten Tag.}{Firma Alpha, Keller, dzień dobry.}
\turn{B}{Guten Tag, hier spricht Piotr Nowak. Kann ich bitte Herrn Meier sprechen?}{Dzień dobry, mówi Piotr Nowak. Czy mogę rozmawiać z panem Meierem?}
\turn{A}{Herr Meier ist leider in einer Besprechung. Kann ich ihm etwas ausrichten?}{Pan Meier jest niestety na naradzie. Czy mogę coś mu przekazać?}
\turn{B}{Ja, bitte. Sagen Sie ihm, dass ich den Termin am Dienstag bestätigen möchte.}{Tak, proszę. Proszę mu powiedzieć, że chcę potwierdzić termin we wtorek.}
\turn{A}{Gern. Wie ist Ihre Nummer?}{Chętnie. Jaki jest Pana numer?}
\end{dialog}

\section{Ćwiczenia}
\begin{cwiczenia}
\exitem{Napisz zwrot początkowy maila do pani Keller (formalnie).}{\de{Sehr geehrte Frau Keller,}}
\exitem{Napisz zwrot początkowy do pana Nowaka (formalnie).}{\de{Sehr geehrter Herr Nowak,}}
\exitem{Napisz zakończenie formalne (DE).}{\de{Mit freundlichen Grüßen}}
\exitem{Napisz zakończenie formalne (CH).}{\de{Freundliche Grüsse}}
\exitem{Poproś o przełożenie spotkania na czwartek.}{\de{Könnten wir den Termin auf Donnerstag verschieben?}}
\exitem{Poproś telefonicznie o połączenie z panem Meierem.}{\de{Kann ich bitte Herrn Meier sprechen?}}
\exitem{Powiedz, że później oddzwonisz.}{\de{Ich rufe später noch einmal an.}}
\exitem{Przedstaw się przez telefon.}{\de{Hier spricht Piotr Nowak.}}
\exitem{Poproś o powtórzenie numeru telefonu.}{\de{Können Sie die Nummer bitte wiederholen?}}
\exitem{Popraw: \de{Witam Pani Keller,}}{\de{Sehr geehrte Frau Keller,}}
\exitem{Przetłumacz: „W załączniku znajdzie Pan dokumenty.”}{\de{Im Anhang finden Sie die Unterlagen.}}
\exitem{Przetłumacz: „Dziękuję za odpowiedź.”}{\de{Vielen Dank für Ihre Rückmeldung.}}
\end{cwiczenia}

\weekend{Napisz trzy e-maile (formalny do klienta, półformalny do zespołu, krótki do kolegi) w sprawie przełożenia
spotkania. Nagraj też 30-sekundową wiadomość na pocztę głosową.}
